\chapter{Tầng 3: Cơ sở Tri thức Đô thị và Hệ thống RAG}

\section{Tổng quan Tầng Tri thức}

Tầng 3 đóng vai trò là \textbf{bộ nhớ pháp lý} của hệ thống STWI --- đảm bảo rằng mọi đề xuất điều phối đều được neo (grounded) vào:
\begin{itemize}[itemsep=3pt]
  \item \textbf{SOP (Standard Operating Procedures):} Quy trình vận hành tiêu chuẩn của Sở GTVT TP.HCM cho từng loại sự cố.
  \item \textbf{Luật GTĐB:} Các điều khoản pháp lý quy định về phân luồng, cấm đường, tốc độ tạm thời.
  \item \textbf{Số liệu mô phỏng cấu trúc:} Kết quả từ Tầng 2 (ma trận lưu lượng, V/C ratio).
\end{itemize}

\section{Kiến trúc Tổng quan Tầng RAG}

\begin{figure}[H]
\centering
\begin{tikzpicture}[node distance=0.5cm and 0.8cm, scale=0.88, transform shape]

  % Ingestion Pipeline (offline)
  \node[font=\small\bfseries, text=tier1col] (offline_label) at (-4.5,3.5)
    {Luồng Nạp dữ liệu (Offline)};

  \node[sysbox=tier1col, text width=2.6cm, minimum height=1.0cm] (sop)
    at (-5.5,2.5) {SOP Sở GTVT\\(PDF/Word)};
  \node[sysbox=tier1col, text width=2.6cm, minimum height=1.0cm] (law)
    at (-3.0,2.5) {Luật GTĐB\\(văn bản pháp lý)};
  \node[sysbox=tier1col, text width=2.6cm, minimum height=1.0cm] (log)
    at (-0.5,2.5) {Log lịch sử\\(case studies)};

  \node[sysbox=tier2col, text width=2.8cm, minimum height=1.0cm] (chunk)
    at (-3.0,1.2) {Chunking\\Sentence-level\\Max 512 tokens};
  \node[sysbox=tier2col, text width=2.8cm, minimum height=1.0cm] (emb)
    at (-3.0,-0.2) {BGE-m3\\Multilingual\\Embedding};
  \node[sysbox=tier3col, text width=3.2cm, minimum height=1.2cm] (vdb)
    at (-3.0,-1.6) {Vector Database\\(Qdrant / Milvus)\\Cosine Similarity};

  \draw[arrow] (sop) -- (chunk);
  \draw[arrow] (law) -- (chunk);
  \draw[arrow] (log) -- (chunk);
  \draw[arrow] (chunk) -- (emb);
  \draw[arrow] (emb) -- (vdb);

  % Query Pipeline (online)
  \node[font=\small\bfseries, text=tier4col] (online_label) at (4.0,3.5)
    {Luồng Truy vấn (Online)};

  \node[sysbox=tier4col, text width=2.8cm, minimum height=1.0cm] (query)
    at (4.0,2.5) {Agent Query\\(ngôn ngữ tự nhiên)};
  \node[sysbox=tier2col, text width=2.8cm, minimum height=1.0cm] (q_emb)
    at (4.0,1.2) {BGE-m3\\Query Embedding};
  \node[sysbox=tier3col, text width=2.8cm, minimum height=1.0cm] (search)
    at (4.0,-0.2) {Semantic Search\\Top-$k$ chunks};
  \node[sysbox=tier4col, text width=2.8cm, minimum height=1.0cm] (ctx)
    at (4.0,-1.6) {Context Assembly\\$+$ Legal Grounding};
  \node[sysbox=tier4col, text width=2.8cm, minimum height=1.2cm,
    fill=alertgreen!50, draw=green!50!black] (ans)
    at (4.0,-3.0) {Agent Response\\với trích dẫn pháp lý};

  \draw[arrow] (query) -- (q_emb);
  \draw[arrow] (q_emb) -- (search);
  \draw[arrow] (search) -- (ctx);
  \draw[arrow] (ctx) -- (ans);

  % Cross-connect
  \draw[dasharrow, bkblue!70] (vdb) -- (search);

  % XiYanSQL branch
  \node[sysbox=tier2col, text width=2.8cm, minimum height=1.0cm] (xsql)
    at (4.0,-4.5) {XiYanSQL\\Text-to-SQL\\In-memory DB};

  \draw[dasharrow, accentorange!80] (ans.south) -- (xsql.north)
    node[midway, right, font=\tiny\itshape]{Nếu cần số liệu};

\end{tikzpicture}
\caption{Kiến trúc tổng quan Tầng RAG của STWI}
\label{fig:rag_arch}
\end{figure}

\section{Cơ sở Dữ liệu Vector}

\subsection{Nguồn Dữ liệu Tích hợp}

\begin{table}[H]
\centering
\caption{Các nguồn dữ liệu được tích hợp vào Vector Database}
\label{tab:knowledge_sources}
\renewcommand{\arraystretch}{1.5}
\begin{tabular}{>{\centering}p{0.8cm} p{4cm} p{7cm}}
\toprule
\textbf{ID} & \textbf{Nguồn dữ liệu} & \textbf{Nội dung và Ví dụ} \\
\midrule
KB-1 & SOP Sở GTVT TP.HCM & Quy trình xử lý: tràn dầu, đổ cây, ngập cục bộ, tai nạn liên hoàn; Phối hợp CSGT, Cứu hỏa, Y tế \\
\rowcolor{rowgray}
KB-2 & Luật GTĐB và Nghị định & Điều kiện cấm đường tạm thời, giới hạn tốc độ đặc biệt, quyền ưu tiên phương tiện, phân luồng hợp pháp \\
KB-3 & Log kịch bản lịch sử & 500+ case studies xử lý sự cố thành công, được dùng làm tiền lệ tham khảo \\
\rowcolor{rowgray}
KB-4 & Bản đồ giao thông & Năng lực đường (capacity), loại đường, số làn, tốc độ thiết kế \\
\bottomrule
\end{tabular}
\end{table}

\subsection{Mô hình Nhúng BGE-m3}

BGE-m3 \cite{chen2023bge} (Beijing Academy of AI General Embedding, Multi-lingual Multi-functionality Multi-granularity) được chọn vì:

\begin{itemize}[itemsep=4pt]
  \item \textbf{Hỗ trợ tiếng Việt tốt:} Được huấn luyện trên corpus đa ngôn ngữ bao gồm tiếng Việt với dấu phụ đầy đủ.
  \item \textbf{Đa chức năng:} Hỗ trợ đồng thời dense retrieval, sparse retrieval và multi-vector retrieval.
  \item \textbf{Độ dài context lớn:} Xử lý được đoạn văn bản dài tới 8.192 tokens.
  \item \textbf{Kích thước embedding:} 1.024 chiều, cân bằng giữa độ chính xác và hiệu năng.
\end{itemize}

\subsection{Chiến lược Chunking}

Chunking là bước quan trọng ảnh hưởng trực tiếp đến chất lượng truy xuất:

\begin{table}[H]
\centering
\caption{Chiến lược chunking theo loại văn bản pháp lý}
\label{tab:chunking}
\renewcommand{\arraystretch}{1.4}
\begin{tabular}{p{3cm} p{3cm} p{6cm}}
\toprule
\textbf{Loại văn bản} & \textbf{Chiến lược} & \textbf{Lý do} \\
\midrule
Điều khoản Luật & Theo điều khoản (Điều 1, Điều 2...) & Mỗi điều khoản là một đơn vị pháp lý độc lập \\
\rowcolor{rowgray}
SOP & Theo bước thực hiện & Mỗi bước SOP cần được tra cứu độc lập \\
Case Studies & Theo sự cố & Mỗi case study là một đơn vị tham chiếu \\
\rowcolor{rowgray}
Văn bản dài & Sliding window 256 tokens, overlap 64 & Đảm bảo không mất ngữ cảnh ở biên chunk \\
\bottomrule
\end{tabular}
\end{table}

\section{Truy xuất Schema-Level với XiYanSQL}

\subsection{Hạn chế của RAG Thuần Văn bản}

Sau khi ADE hoàn thành mô phỏng, kết quả (ma trận V/C ratio, lưu lượng dự báo) được lưu vào In-memory Database có cấu trúc bảng. RAG truyền thống không thể truy vấn chính xác dữ liệu này vì:

\begin{itemize}[itemsep=3pt]
  \item Câu hỏi: ``\textit{V/C ratio tại ngã tư Nguyễn Huệ -- Lê Lợi sau khi đóng đường là bao nhiêu?}''
  \item Semantic search sẽ trả về các đoạn văn bản \textit{mô tả} về V/C ratio, không phải giá trị số thực tế.
  \item Dữ liệu số cần được truy vấn chính xác bằng SQL, không phải xấp xỉ bằng similarity.
\end{itemize}

\subsection{XiYanSQL Pipeline}

\begin{figure}[H]
\centering
\begin{tikzpicture}[node distance=0.5cm and 0.9cm, scale=0.88, transform shape]
  \node[startend] (q) {Câu hỏi NL};
  \node[procnode, right=1.2cm of q, text width=3.5cm] (xi) {XiYanSQL\\Schema-aware\\Text-to-SQL};
  \node[procnode, right=1.2cm of xi, text width=3.5cm] (val) {SQL Validator\\Syntax + Schema\\Check};
  \node[procnode, right=1.2cm of val, text width=3.5cm] (db) {DuckDB\\In-memory\\Execution};
  \node[startend, right=1.2cm of db,
    fill=alertgreen!60, draw=green!50!black, text width=2.8cm] (res)
    {Kết quả số\\chính xác};

  \draw[arrow] (q) -- (xi);
  \draw[arrow] (xi) -- (val);
  \draw[arrow] (val) -- (db);
  \draw[arrow] (db) -- (res);

  % Error path
  \node[procnode, below=1.0cm of val, text width=3cm] (fix) {LLM Self-Correct\\Retry $\leq 3$ lần};
  \draw[arrow] (val.south) -- node[left,font=\tiny]{Lỗi} (fix.north);
  \draw[arrow] (fix.west) -- ++(-0.8,0) |- (xi.south);

  % Schema annotation
  \node[below=1.5cm of xi, font=\tiny, text=codegray, text width=4cm, align=center]
    {Schema cung cấp cho LLM:\\simulation\_results(node\_id,\\timestamp, vc\_ratio,\\avg\_velocity, volume)};
\end{tikzpicture}
\caption{Luồng xử lý XiYanSQL: từ câu hỏi ngôn ngữ tự nhiên đến kết quả SQL}
\label{fig:xiyansql_flow}
\end{figure}

\subsection{Schema In-memory Database}

\begin{lstlisting}[caption={Schema In-memory Database lưu kết quả mô phỏng},label={lst:db_schema}]
-- Schema chinh: Ket qua mo phong Surrogate ADE
CREATE TABLE simulation_results (
    node_id        VARCHAR(50)  NOT NULL,  -- Ma nut giao (e.g., 'NGT_001')
    scenario_id    VARCHAR(50)  NOT NULL,  -- Ma kich ban What-if
    step           INTEGER      NOT NULL,  -- Buoc du bao (1-6, tuong ung 5-30 phut)
    timestamp      TIMESTAMP    NOT NULL,  -- Thoi diem du bao
    vc_ratio       FLOAT        NOT NULL,  -- Volume/Capacity ratio [0, 2.0+]
    avg_velocity   FLOAT        NOT NULL,  -- Van toc trung binh (km/h)
    volume         FLOAT        NOT NULL,  -- Luu luong du bao (PCU/5min)
    congestion_lvl INTEGER      NOT NULL,  -- 0=binh thuong, 1=un u, 2=ket cung
    PRIMARY KEY (node_id, scenario_id, step)
);

-- Schema phu: Ket qua kiem tra CF-VLA
CREATE TABLE cfvla_checks (
    check_id       INTEGER      PRIMARY KEY,
    scenario_id    VARCHAR(50)  NOT NULL,
    iteration      INTEGER      NOT NULL,  -- Vong lap 1, 2, 3
    node_checked   VARCHAR(50)  NOT NULL,
    vc_after       FLOAT        NOT NULL,
    passed         BOOLEAN      NOT NULL,  -- TRUE neu vc_after <= 0.9
    checked_at     TIMESTAMP    DEFAULT CURRENT_TIMESTAMP
);
\end{lstlisting}

\section{RealGen --- Xử lý Kịch bản Cận biên}

\subsection{Thách thức Corner Cases}

Mô hình ADE được huấn luyện trên phân phối lịch sử giao thông. Khi Operator đặt câu hỏi về kịch bản \textit{chưa từng xảy ra} (out-of-distribution), ADE có nguy cơ:
\begin{itemize}[itemsep=3pt]
  \item Extrapolate sai lạc, cho kết quả không thực tế.
  \item Uncertainty cao nhưng không biết cách xử lý.
  \item Dự báo lạc quan quá mức, gây nguy hiểm khi triển khai.
\end{itemize}

\textbf{Ví dụ Corner Case:} \textit{``Tai nạn liên hoàn 5 xe container + mưa lớn 200mm/h + sự kiện đại nhạc hội tại SVĐ Thống Nhất cùng thời điểm.''}

\subsection{Cơ chế RealGen}

\begin{figure}[H]
\centering
\begin{tikzpicture}[node distance=0.5cm and 1.0cm, scale=0.88, transform shape]
  \node[sysbox=tier4col, text width=3cm, minimum height=1.2cm] (q)
    {Kịch bản mới\\(chưa từng có)};

  \node[sysbox=tier2col, right=1.2cm of q, text width=3.2cm, minimum height=1.2cm] (rg)
    {RealGen\\Retrieval Module\\(Similarity Search)};

  \node[sysbox=tier3col, right=1.2cm of rg, text width=3.2cm, minimum height=1.2cm] (hist)
    {Historical\\Corner Cases DB\\(500+ records)};

  \node[sysbox=tier2col, below=1.0cm of rg, text width=3.2cm, minimum height=1.2cm] (aug)
    {Augmentation\\Blend retrieved\\+ current data};

  \node[sysbox=tier1col, right=1.2cm of aug, text width=3.2cm, minimum height=1.2cm] (ade)
    {ADE Surrogate\\Enriched Input\\$\to$ Chính xác hơn};

  \draw[arrow] (q) -- (rg);
  \draw[arrow] (hist) -- (rg);
  \draw[arrow] (rg) -- (aug) node[midway, right, font=\tiny, align=right]{Top-3\\tương tự};
  \draw[arrow] (aug) -- (ade);
  \draw[arrow] (q.south) -- ++(0,-1.0) -| (aug.west);
\end{tikzpicture}
\caption{Cơ chế RealGen xử lý kịch bản cận biên}
\label{fig:realgen}
\end{figure}

\subsection{Chỉ số Tương đồng Kịch bản}

RealGen tính độ tương đồng giữa kịch bản mới $\mathbf{c}_{\text{new}}$ và kịch bản lịch sử $\mathbf{c}_{\text{hist},j}$ theo công thức:

\begin{equation}
  \text{sim}(\mathbf{c}_{\text{new}}, \mathbf{c}_{\text{hist},j}) = \alpha \cdot \cos(\mathbf{e}_{\text{new}}, \mathbf{e}_j) + (1-\alpha) \cdot \exp\!\left(-\frac{\|\mathbf{c}_{\text{new}} - \mathbf{c}_j\|_2}{\tau}\right)
  \label{eq:realgen_sim}
\end{equation}

với $\mathbf{e}$ là text embedding của mô tả kịch bản, $\alpha = 0{,}6$, $\tau = 2{,}0$.

\section{Chính sách Không Ảo giác (No-Hallucination Policy)}

\begin{tcolorbox}[dangerbox]
\textbf{Quy tắc bất biến (QC-7):} Mọi response từ Legal Agent \textbf{phải kèm theo trường \texttt{legal\_grounding}} với định dạng: \texttt{"[Tên nguồn] Điều X, Khoản Y: [Trích dẫn nguyên văn]"}. Response thiếu trường này sẽ bị từ chối bởi Evaluation Agent.
\end{tcolorbox}

Cơ chế đảm bảo No-Hallucination:
\begin{enumerate}[leftmargin=1.5cm, itemsep=4pt]
  \item \textbf{Source Attribution:} Mỗi chunk trong Vector DB gắn metadata nguồn gốc (tên văn bản, điều khoản, ngày ban hành).
  \item \textbf{Citation Injection:} System prompt bắt buộc LLM trích dẫn chính xác nguồn được retrieve.
  \item \textbf{Hallucination Detector:} Evaluation Agent kiểm tra xem \texttt{legal\_grounding} có khớp với nội dung thực sự trong Vector DB hay không bằng cosine similarity $> 0{,}85$.
  \item \textbf{Fallback Label:} Nếu không tìm được căn cứ phù hợp, Agent tự động gắn nhãn \texttt{[CANH BAO -- CHUA KIEM CHUNG]} thay vì bịa đặt.
\end{enumerate}
